% GENERATED FILE. Edit canonical lessons, then run npm run generate:books.
% canonical-source-hash: fnv1a64:390b3ad924a641ed
% canonical-lessons: ZH-C414-chengwuyuan, ZH-C414-kaochuang, ZH-C414-jingjicang, ZH-C414-yunji, ZH-C414-anquandai

\chapter{In the Cabin}
\label{ch:zh-in-the-cabin}
\hlchaptermodality{414}

\begin{chapteropening}
\textbf{By the end of this chapter:} \emph{I can talk about the cabin crew, a window seat, economy class, feeling sick and the seat belt.}

Complete the last lesson of chapter 414: I can talk about the cabin crew, a window seat, economy class, feeling sick and the seat belt.
\end{chapteropening}

\section[chéngwùyuán]{chéngwùyuán — flight attendant; cabin crew}

\label{lesson:ZH-C414-chengwuyuan}



\begin{quote}
Before the new one: say the Mandarin for to come down onto the runway, then the Mandarin for as fast as flying.
\end{quote}

\subsection*{You'll want to know: chéngwùyuán}
\emph{chéngwùyuán} — \textquotedblleft{}flight attendant; cabin crew\textquotedblright{}.

\begin{quote}
\emph{chéngwùyuán zhǎohuí le wǒ de shǒujī} — The flight attendant found my phone for me.
\end{quote}

Its characters come later; for now it is a word for the ear.

\begin{cousinweb}
\emph{chéngwùyuán} is built from \emph{chéng} 'to ride', \emph{wù} 'a task' and \emph{yuán} 'a member of staff', all inherited from Middle Chinese, \emph{zying}, \emph{mjuH} and \emph{hjwen} (Guangyun, in Baxter's transcription), which glosses the second 'business, affairs'.
\end{cousinweb}

\begin{grammarlens}[title={What you've built}]
One more word for getting around.
\end{grammarlens}

\subsection*{Your turn}
\begin{itemize}
\raggedright
  \item \emph{Say it:} \emph{chéngwùyuán}
  \item \emph{Say it:} \emph{chéngwùyuán}, once more
  \item \emph{Say it:} \emph{chéngwùyuán zhǎohuí le wǒ de shǒujī} — The flight attendant found my phone for me.
  \item \emph{Recall:} say the Mandarin for to come down onto the runway, then the Mandarin for as fast as flying, then say \emph{chéngwùyuán} again
\end{itemize}

\begin{tcolorbox}[breakable,colback=teal!4,colframe=teal!35!black,title={Before you move on}]
What does chéngwùyuán mean? (\textquotedblleft{}Flight attendant; cabin crew\textquotedblright{}.) Say it once more.
\end{tcolorbox}
\section[kàochuāng]{kàochuāng — next to the window}

\label{lesson:ZH-C414-kaochuang}



\begin{quote}
Before the new one: say the Mandarin for as fast as flying, then the Mandarin for a flight attendant.
\end{quote}

\subsection*{You'll want to know: kàochuāng}
\emph{kàochuāng} — \textquotedblleft{}next to the window\textquotedblright{}.

\begin{quote}
\emph{wǒ xiǎng yào yí ge kàochuāng de zuòwèi} — I'd like a window seat.
\end{quote}

Its characters come later; for now it is a word for the ear.

\begin{cousinweb}
\emph{kàochuāng} is built from \emph{kào} 'to be next to' and \emph{chuāng} 'a window', both inherited from Middle Chinese, \emph{khawH} and \emph{tsrhaewng} (Guangyun, in Baxter's transcription), which glosses the second, after an older dictionary, as an opening through a wall.
\end{cousinweb}

\begin{grammarlens}[title={What you've built}]
One more word for getting around.
\end{grammarlens}

\subsection*{Your turn}
\begin{itemize}
\raggedright
  \item \emph{Say it:} \emph{kàochuāng}
  \item \emph{Say it:} \emph{kàochuāng}, once more
  \item \emph{Say it:} \emph{wǒ xiǎng yào yí ge kàochuāng de zuòwèi} — I'd like a window seat.
  \item \emph{Recall:} say the Mandarin for as fast as flying, then the Mandarin for a flight attendant, then say \emph{kàochuāng} again
\end{itemize}

\begin{tcolorbox}[breakable,colback=teal!4,colframe=teal!35!black,title={Before you move on}]
What does kàochuāng mean? (\textquotedblleft{}Next to the window\textquotedblright{}.) Say it once more.
\end{tcolorbox}
\section[jīngjìcāng]{jīngjìcāng — economy class}

\label{lesson:ZH-C414-jingjicang}



\begin{quote}
Before the new one: say the Mandarin for a flight attendant, then the Mandarin for next to the window.
\end{quote}

\subsection*{You'll want to know: jīngjìcāng}
\emph{jīngjìcāng} — \textquotedblleft{}economy class\textquotedblright{}.

\begin{quote}
\emph{wǒmen zuò jīngjìcāng qù Běijīng} — We're flying to Beijing in economy.
\end{quote}

Its characters come later; for now it is a word for the ear.

\begin{cousinweb}
\emph{jīngjìcāng} is built from \emph{jīngjì} 'economy' and \emph{cāng} 'a cabin'. The first two parts are inherited from Middle Chinese, \emph{keng} and \emph{tsejH} (Guangyun, in Baxter's transcription), glossed 'constant' and 'to cross a river'; no dictionary we use records an older form of the third, so its history is not settled.
\end{cousinweb}

\begin{grammarlens}[title={What you've built}]
One more word for getting around.
\end{grammarlens}

\subsection*{Your turn}
\begin{itemize}
\raggedright
  \item \emph{Say it:} \emph{jīngjìcāng}
  \item \emph{Say it:} \emph{jīngjìcāng}, once more
  \item \emph{Say it:} \emph{wǒmen zuò jīngjìcāng qù Běijīng} — We're flying to Beijing in economy.
  \item \emph{Recall:} say the Mandarin for a flight attendant, then the Mandarin for next to the window, then say \emph{jīngjìcāng} again
\end{itemize}

\begin{tcolorbox}[breakable,colback=teal!4,colframe=teal!35!black,title={Before you move on}]
What does jīngjìcāng mean? (\textquotedblleft{}Economy class\textquotedblright{}.) Say it once more.
\end{tcolorbox}
\section[yùnjī]{yùnjī — to get airsick}

\label{lesson:ZH-C414-yunji}



\begin{quote}
Before the new one: say the Mandarin for next to the window, then the Mandarin for economy class.
\end{quote}

\subsection*{You'll want to know: yùnjī}
\emph{yùnjī} — \textquotedblleft{}to get airsick\textquotedblright{}.

\begin{quote}
\emph{wǒ zuò fēijī de shíhou huì yùnjī} — I get airsick when I fly.
\end{quote}

Its characters come later; for now it is a word for the ear.

\begin{cousinweb}
\emph{yùnjī} is built from \emph{yùn} 'dizzy' and \zh{机} 'a plane'. The second is inherited from Middle Chinese \emph{kj+j} (Guangyun, in Baxter's transcription); the dictionary glosses the first, \emph{hjunH} there, as a halo round the sun or moon, so how it came to mean dizziness is not settled.
\end{cousinweb}

\begin{grammarlens}[title={What you've built}]
One more word for getting around.
\end{grammarlens}

\subsection*{Your turn}
\begin{itemize}
\raggedright
  \item \emph{Say it:} \emph{yùnjī}
  \item \emph{Say it:} \emph{yùnjī}, once more
  \item \emph{Say it:} \emph{wǒ zuò fēijī de shíhou huì yùnjī} — I get airsick when I fly.
  \item \emph{Recall:} say the Mandarin for next to the window, then the Mandarin for economy class, then say \emph{yùnjī} again
\end{itemize}

\begin{tcolorbox}[breakable,colback=teal!4,colframe=teal!35!black,title={Before you move on}]
What does yùnjī mean? (\textquotedblleft{}To get airsick\textquotedblright{}.) Say it once more.
\end{tcolorbox}
\section[ānquándài]{ānquándài — seat belt}

\label{lesson:ZH-C414-anquandai}



\begin{quote}
Before the new one: say the Mandarin for economy class, then the Mandarin for to get airsick.
\end{quote}

\subsection*{You'll want to know: ānquándài}
\emph{ānquándài} — \textquotedblleft{}seat belt\textquotedblright{}.

\begin{quote}
\emph{měi ge zuòwèi dōu yǒu ānquándài} — Every seat on the plane has a seat belt.
\end{quote}

Its characters come later; for now it is a word for the ear.

\begin{cousinweb}
\emph{ānquándài} is built from \emph{ānquán} 'safe' and \emph{dài} 'a belt', all inherited from Middle Chinese, \emph{'an}, \emph{dzjwen} and \emph{tajH} (Guangyun, in Baxter's transcription), which glosses the last as a sash worn with clothes.
\end{cousinweb}

\begin{grammarlens}[title={What you've built}]
One more word for getting around.
\end{grammarlens}

\subsection*{Your turn}
\begin{itemize}
\raggedright
  \item \emph{Say it:} \emph{ānquándài}
  \item \emph{Say it:} \emph{ānquándài}, once more
  \item \emph{Say it:} \emph{měi ge zuòwèi dōu yǒu ānquándài} — Every seat on the plane has a seat belt.
  \item \emph{Recall:} say the Mandarin for economy class, then the Mandarin for to get airsick, then say \emph{ānquándài} again
\end{itemize}

\begin{tcolorbox}[breakable,colback=teal!4,colframe=teal!35!black,title={Before you move on}]
What does ānquándài mean? (\textquotedblleft{}Seat belt\textquotedblright{}.) Say it once more.
\end{tcolorbox}
